\section{Density lower bounds and mixed products}
\label{sec:bt-density-lower-bounds}

Let $\Omega\subset\mathbb C^n$ be open, with $n\ge1$, and put $\eta=\ddc|z|^2=\beta/\pi$. Thus $\eta^n=n!\pi^{-n}\,\mathrm dV_{2n}$ by \eqref{eq:bt-determinant-normalization}. A lower bound for a Monge--Amp\`ere density can be tested against smooth functions touching the potential from above. A matrix formula then turns these tests into linear differential inequalities, which survive convolution. This permits the smooth mixed determinant inequality to pass to locally bounded potentials.

Here a real $C^2$ function $q$ touches a real function $u$ from above at $x$ if $q(x)=u(x)$ and $q\ge u$ on a neighborhood of $x$.

\begin{lemma}\label{lem:determinant-linearization-mixed}
For a positive semidefinite Hermitian $n\times n$ matrix $A$,
\[
 (\det A)^{1/n}=\inf_{B>0,\,\det B=1}\frac{\operatorname{tr}(BA)}n.
\]
In the infimum $B>0$ means positive definite Hermitian. For positive semidefinite $A,C$, define $M_k(A,C)$ by
\[
 \det(tA+sC)=\sum_{k=0}^n\binom nk M_k(A,C)t^ks^{n-k}.
\]
For each integer $1\le k\le n-1$,
\[
 M_k(A,C)\ge(\det A)^{k/n}(\det C)^{(n-k)/n}.
\]
The endpoint coefficients are $M_n(A,C)=\det A$ and $M_0(A,C)=\det C$.
\end{lemma}
\begin{proof}
If $A>0$, apply the arithmetic--geometric mean inequality to the eigenvalues of $B^{1/2}AB^{1/2}$. Equality is attained at $B=(\det A)^{1/n}A^{-1}$. If $A$ is singular, use this choice for $A+\varepsilon I$; its trace against $A$ is at most $n\det(A+\varepsilon I)^{1/n}$ and tends to zero. All the traces are nonnegative, proving the first assertion.

Suppose first that $A,C>0$ and let $\lambda_1,\ldots,\lambda_n$ be the eigenvalues of $A^{-1/2}CA^{-1/2}$. Coefficient comparison gives
\[
 M_k(A,C)=\frac{\det A}{\binom nk}
 \sum_{|J|=n-k}\prod_{j\in J}\lambda_j.
\]
Each eigenvalue occurs in $\binom{n-1}{n-k-1}$ of these products. Their geometric mean is therefore $(\lambda_1\cdots\lambda_n)^{(n-k)/n}$. The arithmetic--geometric mean inequality gives the stated bound since $\prod_j\lambda_j=\det C/\det A$. Replace both matrices by their sums with $\varepsilon I$ and let $\varepsilon\downarrow0$ to include the semidefinite case.
\end{proof}

\begin{lemma}\label{lem:constant-elliptic-upper-tests}
Let $d\ge1$ be an integer and $L=\sum_{a,b=1}^d a_{ab}\partial_a\partial_b$, where $(a_{ab})$ is a real symmetric positive definite constant matrix. Let $w$ be real, locally bounded and upper semicontinuous on an open $U\subset\mathbb R^d$, and let $F:U\to\mathbb R$ be continuous. Suppose every $C^2$ function $q$ touching $w$ from above at $x$ satisfies $Lq(x)\ge F(x)$. Then $Lw\ge F$ distributionally.
\end{lemma}
\begin{proof}
A real linear change of variables makes $L$ the ordinary Laplacian. If $F=0$, the upper-test condition implies harmonic comparison on every relatively compact ball $B(y,r)$. Indeed, let $h$ be harmonic on that ball, continuous on its closure, and at least $w$ on the boundary. If $w-h$ is positive somewhere inside, choose $\varepsilon>0$ so small that
$w-h-\varepsilon(r^2-|x-y|^2)$ still has a positive value there. Its positive maximum is attained inside, since its boundary upper limits are nonpositive. Adding this maximum to $h+\varepsilon(r^2-|x-y|^2)$ gives an upper test with strictly negative Laplacian, a contradiction. Decreasing continuous majorants of $w|_{\partial B}$ and their Poisson extensions gives the sphere submean inequality, as in \cref{thm:subharmonic-comparison}. Consequently \cref{thm:subharmonic-distribution} gives $\Delta w\ge0$.

For general $F$, work on a smaller ball and choose smooth $F_j\le F$ converging uniformly to $F$ there. Such functions result from convolution on a larger neighborhood followed by subtraction of its uniform approximation error. Multiply $F_j$ by a smooth cutoff equal to one near the smaller ball. Convolution with the fundamental solution in \cref{prop:laplacian-fundamental-solution} gives a smooth $h_j$ with $\Delta h_j=F_j$ on that ball: derivatives may be placed on the compactly supported smooth factor. Every upper test for $w-h_j$ has nonnegative Laplacian, so the first case gives $\Delta w\ge F_j$ distributionally. Uniform convergence gives $\Delta w\ge F$. Returning to the original coordinates proves the lemma.
\end{proof}

For a real $C^2$ function $q$ write $H(q)=(q_{j\bar k})$. For a positive definite Hermitian $B$ with $\det B=1$, put $L_Bq=\operatorname{tr}(BH(q))$. This is a constant real elliptic operator: diagonalizing $B$ by a unitary coordinate change expresses it as a positive weighted sum of the planar Laplacians, each with factor $1/4$.

\begin{proposition}\label{prop:bounded-monge-ampere-linear-density}
Let $u$ be locally bounded and plurisubharmonic on $\Omega$, and let $f\ge0$ be continuous. The following conditions are equivalent:
\begin{enumerate}
\item $(\ddc u)^n\ge f\eta^n$ as measures;
\item every real $C^2$ function $q$ touching $u$ from above at $x$ satisfies $\det H(q)(x)\ge f(x)$;
\item $L_Bu\ge n f^{1/n}$ distributionally for every positive definite Hermitian $B$ with $\det B=1$.
\end{enumerate}
\end{proposition}
\begin{proof}
The complex Hessian of any upper test $q$ is positive semidefinite at the contact point. Restrict to a complex line through that point and use the submean inequality for $u$; the circle mean of $q$ is at least its center value. The second-order expansion gives nonnegativity of $H(q)$ on every complex direction.

Assume (1), and normalize a test by $q(x)=u(x)$ and $q\ge u$ near $x$. If $\det H(q)(x)<f(x)$, choose $\delta>0$ so small that $q_\delta=q+\delta|z-x|^2$ has positive definite complex Hessian at $x$ and still has determinant below $f(x)$. On a sufficiently small closed ball $\overline{B(x,r)}$, it is strictly plurisubharmonic and $\det H(q_\delta)<f$. Set $b=q_\delta-\delta r^2/2$. On the sphere $b>u$, whereas $b(x)<u(x)$. The Borel set $S=\{b<u\}\cap B(x,r)$ has positive Lebesgue measure. Otherwise $u\le b$ almost everywhere on the ball, and uniqueness of the plurisubharmonic representative in \cref{thm:psh-distribution-representative}, applied to $\max\{u,b\}$, would imply $u\le b$ everywhere. Bounded comparison in \cref{thm:bt-comparison} now gives
\[
 \int_S f\eta^n\le\int_S(\ddc u)^n
 \le\int_S(\ddc b)^n
 =\int_S\det H(q_\delta)\eta^n,
\]
contradicting the strict density inequality. The measures here are finite since the ball is relatively compact in the neighborhood where the potentials are bounded. Thus (2) holds.

Assume (2). The first formula of \cref{lem:determinant-linearization-mixed} gives $L_Bq(x)\ge n f(x)^{1/n}$ for every upper test and each fixed admissible $B$. Apply \cref{lem:constant-elliptic-upper-tests} with the continuous function $n f^{1/n}$ to obtain (3).

Finally assume (3). Choose the nonnegative radial regularizing kernels used in \cref{thm:psh-distribution-representative}, so that $u_\varepsilon=u*\rho_\varepsilon\downarrow u$ on smaller open sets. Convolving each distributional inequality gives
\[
 \operatorname{tr}(BH(u_\varepsilon))\ge n(f^{1/n}*\rho_\varepsilon)
 \qquad(B>0,\ \det B=1).
\]
These are pointwise inequalities between smooth functions. The matrix infimum formula therefore gives $\det H(u_\varepsilon)\ge(f^{1/n}*\rho_\varepsilon)^n$. The right side tends locally uniformly to $f$. The determinant normalization and the weak decreasing continuity of \cref{thm:bt-decreasing-continuity} give (1). This argument includes vanishing densities and singular Hessians.
\end{proof}
\index{Monge--Amp\`ere measure!density lower bound}

\begin{corollary}\label{cor:bounded-mixed-continuous-density}
Let $u,v$ be locally bounded plurisubharmonic functions on $\Omega$, and let $f,g\ge0$ be continuous with $(\ddc u)^n\ge f\eta^n$ and $(\ddc v)^n\ge g\eta^n$. For each integer $1\le k\le n-1$,
\[
 (\ddc u)^k\wedge(\ddc v)^{n-k}
 \ge f^{k/n}g^{(n-k)/n}\eta^n.
\]
For $k=0,n$ the corresponding assertions are the given pure-measure bounds.
\end{corollary}
\begin{proof}
Use the same radial kernels for both potentials. The proof of \cref{prop:bounded-monge-ampere-linear-density} gives
\[
 \det H(u_\varepsilon)\ge(f^{1/n}*\rho_\varepsilon)^n,
 \qquad
 \det H(v_\varepsilon)\ge(g^{1/n}*\rho_\varepsilon)^n.
\]
Expand $(t\ddc u_\varepsilon+s\ddc v_\varepsilon)^n$ and its smooth determinant density. The coefficient of $\binom nk t^ks^{n-k}$ is $M_k(H(u_\varepsilon),H(v_\varepsilon))\eta^n$. The second assertion of \cref{lem:determinant-linearization-mixed} bounds this below by
$(f^{1/n}*\rho_\varepsilon)^k(g^{1/n}*\rho_\varepsilon)^{n-k}\eta^n$.
The continuous density converges locally uniformly to $f^{k/n}g^{(n-k)/n}$. Mixed decreasing continuity in \cref{thm:bt-decreasing-continuity} passes the measure inequality to the limit.
\end{proof}
\index{Bedford--Taylor product!mixed density inequality}
